
\documentclass{article} 
\usepackage{collas2026_conference,times}
\usepackage{easyReview}

\usepackage[utf8]{inputenc} 
\usepackage[T1]{fontenc}    
\usepackage{hyperref}      
\usepackage{url}            
\usepackage{booktabs}       
\usepackage{amsfonts}      
\usepackage{nicefrac}     
\usepackage{microtype}     
\usepackage{xcolor}         
\usepackage{multirow}
\usepackage{wrapfig,lipsum,booktabs}
\usepackage{booktabs}
\usepackage{graphicx}
\usepackage{amsmath}
\usepackage{subcaption}
\usepackage{placeins}
\usepackage{pifont}
\newcommand{\cmark}{\ding{51}}
\newcommand{\xmark}{\ding{55}}

\newcommand{\elena}[1]{{\color{magenta}{[Elena: #1]}}}
\newcommand{\tibi}[1]{{\color{blue}{[Tibi: #1]}}}


\input{math_commands.tex}

\usepackage{hyperref}
\hypersetup{
    colorlinks=true,
    linkcolor=red,
    filecolor=magenta,
    urlcolor=blue,
    citecolor=purple,
    pdftitle={The Formal Theory of Reversible Behavioral Learning},
    pdfpagemode=FullScreen,
    }


\title{The Formal Theory of Reversible Behavioral Learning: \\
Path Functionals and Behavioral Geometry}


\author{Pardhu Sri Rushi Varma Konduru \\
    Malla Reddy University \\
    Hyderabad, Telangana, India \\
    \texttt{pardhuvarma.cs@gmail.com}
}


\newcommand{\fix}{\marginpar{FIX}}
\newcommand{\new}{\marginpar{NEW}}

%\collasfinalcopy % Uncomment for camera-ready version, but NOT for submission.

\preprintcopy 

\begin{document}


\maketitle


\begin{abstract}
Prior work established structural irreversibility as a fundamental property of weight-based neural 
adaptation and introduced Runtime Low-Rank Adaptive Environments (RLAE), a paradigm in which 
adaptive behavior is encoded in a separable parameter set $\phi$, admitting deterministic rollback 
via an explicit unload operator $\mathcal{K}$, with recoverability validated empirically across 
model scales. That framework established \textit{that} behavioral change can be structurally 
reversed---but left open \textit{how} behavioral change evolves within the reversible setting.

This paper develops a geometric and variational theory of behavioral evolution within RLAE. We 
model adaptation as a continuous trajectory in behavioral parameter space $\Phi$, introduce 
Behavioral Path Functionals (BPF) as cost functionals over such trajectories, and instantiate 
them via the Informational Geodesic Path Functional (IGPF), which assigns cost through 
information-theoretic divergence between successive behavioral states. Extending SVAR, we show 
that local behavioral sensitivity admits a quadratic approximation inducing a Riemannian metric 
over $\Phi$, from which Geodesic Behavioral Distance (GBD) is defined as the minimal-cost 
transformation between behavioral states. Variational analysis of the IGPF yields the central 
result: optimal behavioral evolution follows geodesics, establishing a formal connection to natural 
gradient methods.

These results extend RLAE from a structural recoverability guarantee to a geometric theory of 
behavioral dynamics, establishing adaptation as a path-dependent process governed by information 
geometry over $\Phi$.
\end{abstract}


\section{Introduction}\label{section-introduction}

Modern neural systems increasingly rely on post-training adaptation to incorporate new tasks,
align behavior, and operate under changing objectives. A central question in this setting is
whether adaptive changes can be governed, audited, and reversed without compromising the
integrity of the underlying model. In prior work~\cite{konduru2026structurallimitationsweightbasedneural}, we examined adaptation
from a structural perspective and identified a fundamental limitation of weight-based approaches:
modifications applied directly to shared model parameters $\theta$ introduce behavioral changes
that are entangled with core model representations and cannot, in general, be deterministically
undone. We formalized this as \textit{structural irreversibility} and demonstrated empirically
that rollback under direct weight mutation requires either explicit parameter checkpoints or
retraining, with no guarantee of behavioral equivalence.

To address this limitation, we introduced \textit{Runtime Low-Rank Adaptive Environments}
(RLAE) in prior work, a paradigm in which adaptive behavior is represented exclusively in a separable
behavioral parameter set $\phi$, leaving the core parameter set $\theta$ frozen throughout
adaptation. An explicit unload operator $\mathcal{K}$ removes the behavioral component
entirely, restoring the model to its baseline identity without approximation or retraining.
This establishes reversibility as a structural property of the adaptation mechanism itself,
validated empirically via the Recoverability Factor (RF), Identity Leakage Score (ILS), and
Structural Variance Analysis for Robustness (SVAR) across multiple model scales.

Earlier in prior work, we introduced RLAE as a structural framework for
reversible behavioral adaptation, focusing on identity preservation and deterministic rollback
under parameter isolation. The scope of that work was centered on the existence and
reversibility of behavioral states, establishing recoverability as a first-class property of
adaptive systems.

Building on this foundation, a natural next step is to examine not only whether behavioral
states can be reached and reversed, but how transitions between such states are structured.
In particular, the RLAE formulation treats adaptation as a sequence of discrete updates to
$\phi$, without explicitly modeling continuity, path dependence, or the geometry of
behavioral change over $\Phi$.

This motivates a shift from state-level reasoning to a path-based formalization of
adaptation, in which behavioral evolution is treated as a continuous transformation with
associated cost, structure, and optimality conditions. In this work, we develop this
perspective through the introduction of Behavioral Path Functionals (BPF), which assign 
a scalar cost to trajectories in $\Phi$, enabling adaptation to be formulated as an
optimization problem over paths rather than over parameter values directly. We instantiate
this general framework through the \textit{Informational Geodesic Path Functional} (IGPF),
in which trajectory cost is induced by information-theoretic divergence between successive
behavioral states. This extends the KL-based divergence metrics introduced in prior work
from static diagnostic tools into a variational structure governing behavioral transitions.

Geometric structure over $\Phi$ arises naturally from SVAR. The quadratic approximation of
local behavioral sensitivity under bounded perturbations of $\phi$ defines a Riemannian
metric over the behavioral parameter space, without modifying $\theta$ or the unload
semantics of RLAE. This metric gives rise to \textit{Geodesic Behavioral Distance} (GBD),
the minimal-cost transformation between behavioral states under the induced geometry.
Variational analysis of the IGPF then yields the central result of this paper: optimal
behavioral evolution follows geodesics under this metric, establishing a formal and
principled connection to natural gradient methods.

The contributions of this paper are as follows:

\begin{itemize}
    \item We introduce \textit{Behavioral Path Functionals} (BPF) as a general variational
    framework for modeling adaptation as a path-dependent process over the behavioral
    parameter space $\Phi$.

    \item We define the \textit{Informational Geodesic Path Functional} (IGPF), which
    grounds trajectory cost in information-theoretic divergence, extending the divergence
    metrics of prior work into a variational setting.

    \item We derive a Riemannian metric over $\Phi$ from SVAR, establishing a natural
    geometric structure over the behavioral parameter space that is consistent with the
    frozen-core constraint of RLAE.

    \item We formalize \textit{Geodesic Behavioral Distance} (GBD) and establish through
    variational analysis that optimal behavioral evolution follows geodesics, connecting
    reversible adaptation to natural gradient methods.
\end{itemize}

Together, these results extend RLAE from a structural recoverability guarantee to a complete
geometric theory of behavioral dynamics. Adaptation within the reversible setting is
reframed as a path-dependent process governed by information geometry over $\Phi$, providing
a principled foundation for analyzing the efficiency, stability, and controllability of
neural adaptation in modular and long-lived systems.

\newpage

\bibliography{collas2026_conference}
\bibliographystyle{collas2026_conference}

\newpage

\appendix

\section*{Appendix}

\section{Additional Theoretical Details}
\label{appendix:proofs}

For completeness, we restate the progress bound used in the main text.

\paragraph{Progress bound under projected updates.}
Let
\begin{equation}
\theta^+ = \theta - \eta P_S \nabla \mathcal{J}_t(\theta),
\label{eq:appendix_theta_plus}
\end{equation}
where $P_S$ is the orthogonal projector associated with the fixed trainable subset $S$. Assume that $\mathcal{J}_t$ is $L$-smooth on a neighborhood containing both $\theta$ and $\theta^+$. Then
\begin{equation}
\mathcal{J}_t(\theta^+)
\le
\mathcal{J}_t(\theta)
-
\eta \|P_S \nabla \mathcal{J}_t(\theta)\|^2
+
\frac{L\eta^2}{2}\|P_S \nabla \mathcal{J}_t(\theta)\|^2.
\label{eq:appendix_descent_1}
\end{equation}
Equivalently,
\begin{equation}
\mathcal{J}_t(\theta^+)
\le
\mathcal{J}_t(\theta)
-
\eta \Bigl(1-\frac{L\eta}{2}\Bigr)\|P_S \nabla \mathcal{J}_t(\theta)\|^2.
\label{eq:appendix_descent_2}
\end{equation}
In particular, if $0 < \eta \le 1/L$, then
\begin{equation}
\mathcal{J}_t(\theta^+)
\le
\mathcal{J}_t(\theta)
-
\frac{\eta}{2}\|P_S \nabla \mathcal{J}_t(\theta)\|^2.
\label{eq:appendix_descent_3}
\end{equation}

\paragraph{Proof.}
Since $\mathcal{J}_t$ is $L$-smooth, for any $x,y$ in the neighborhood of interest,
\begin{equation}
\mathcal{J}_t(y)
\le
\mathcal{J}_t(x)
+
\langle \nabla \mathcal{J}_t(x), y-x \rangle
+
\frac{L}{2}\|y-x\|^2.
\label{eq:appendix_smoothness}
\end{equation}
Apply \ref{eq:appendix_smoothness} with $x=\theta$ and $y=\theta^+$. Using \ref{eq:appendix_theta_plus},
\begin{equation}
\theta^+ - \theta = -\eta P_S \nabla \mathcal{J}_t(\theta).
\label{eq:appendix_step_diff}
\end{equation}
Substituting into \ref{eq:appendix_smoothness} gives
\begin{align}
\mathcal{J}_t(\theta^+)
&\le
\mathcal{J}_t(\theta)
+
\left\langle \nabla \mathcal{J}_t(\theta), \theta^+ - \theta \right\rangle
+
\frac{L}{2}\|\theta^+ - \theta\|^2 \\
&=
\mathcal{J}_t(\theta)
-
\eta \left\langle \nabla \mathcal{J}_t(\theta), P_S \nabla \mathcal{J}_t(\theta) \right\rangle
+
\frac{L\eta^2}{2}\|P_S \nabla \mathcal{J}_t(\theta)\|^2.
\label{eq:appendix_before_projector_identity}
\end{align}

Now use that $P_S$ is an orthogonal projector. Hence it is symmetric and idempotent:
\begin{equation}
P_S^\top = P_S,
\qquad
P_S^2 = P_S.
\label{eq:appendix_projector_properties}
\end{equation}
Therefore,
\begin{align}
\left\langle \nabla \mathcal{J}_t(\theta), P_S \nabla \mathcal{J}_t(\theta) \right\rangle
&=
\nabla \mathcal{J}_t(\theta)^\top P_S \nabla \mathcal{J}_t(\theta) \\
&=
\nabla \mathcal{J}_t(\theta)^\top P_S^\top P_S \nabla \mathcal{J}_t(\theta) \\
&=
\|P_S \nabla \mathcal{J}_t(\theta)\|^2.
\label{eq:appendix_projector_identity}
\end{align}
Substituting \ref{eq:appendix_projector_identity} into \ref{eq:appendix_before_projector_identity} yields \ref{eq:appendix_descent_1}. Rearranging gives \ref{eq:appendix_descent_2}. Finally, if $0<\eta\le 1/L$, then
\begin{equation}
1-\frac{L\eta}{2}\ge \frac{1}{2},
\end{equation}
which implies \ref{eq:appendix_descent_3}. \hfill $\square$

\newpage

\section{Notation Summary}
In Tab.~\ref{tab:notation}, we provide an overview of the mathematical notation.

\begin{table}[!h]
\centering
\small
\caption{Notation used in the theoretical analysis}
\begin{tabular}{p{0.20\linewidth} p{0.72\linewidth}}
\toprule
\textbf{Symbol} & \textbf{Meaning} \\
\midrule
$\mathcal{T}_1, \dots, \mathcal{T}_T$ & Sequence of tasks observed in continual learning. \\

$T$ & Number of tasks in the continual learning sequence. \\

$\theta \in \mathbb{R}^d$ & Full model parameter vector. \\

$d$ & Total number of model parameters. \\

$t$ & Task index. \\

$m$ & Optimization-step index within task $t$. \\

$\theta_t$ & Parameters after completing training on task $t$. \\

$\theta_t^{(m)}$ & Parameters after $m$ optimization steps on task $t$. \\

$\mathcal{L}_t(\theta)$ & Current-task loss at task $t$. \\

$\Omega_t(\theta;\theta_{1:t-1})$ & Method-specific term that preserves previously acquired knowledge. \\

$\lambda$ & Weight of the preservation term. \\

$\mathcal{J}_t(\theta)$ & Full task-$t$ objective, defined as $\mathcal{L}_t(\theta) + \lambda \Omega_t(\theta;\theta_{1:t-1})$. \\

$S \subseteq \{1,\dots,d\}$ & Fixed subset of trainable parameter coordinates defining an adaptation regime. \\

$P_S$ & Orthogonal projector that keeps coordinates in $S$ and zeros out all others. \\

$\eta$ & Step size used by the first-order optimizer. \\

$g_t^{(S)}(\theta)$ & Projected current-task update signal, defined as $P_S \nabla \mathcal{L}_t(\theta)$. \\

$r_t^{(S)}(\theta)$ & Projected preservation update signal, defined as $P_S \nabla \Omega_t(\theta;\theta_{1:t-1})$. \\

$\Gamma_t(S;\theta)$ & Interaction term between projected current-task and preservation signals, defined as $\langle g_t^{(S)}(\theta), r_t^{(S)}(\theta)\rangle$. \\

$P_S \nabla \mathcal{J}_t(\theta)$ & Projected gradient of the full task-$t$ objective. \\

$L$ & Smoothness constant of $\mathcal{J}_t$. \\

$\theta^+$ & One-step projected update from $\theta$, defined in \ref{eq:appendix_theta_plus}. \\

$\|\cdot\|$ & Euclidean norm. \\

$\langle \cdot,\cdot \rangle$ & Standard Euclidean inner product. \\
\bottomrule
\end{tabular}
\label{tab:notation}
\end{table}

\section{Additional experimental values}

Tables~\ref{tab:cifar100}, \ref{tab:mnist}, \ref{tab:fashion}, \ref{tab:qmnist}, and \ref{tab:kmnist} report the numerical values underlying the results shown in Figures \ref{fig:main_results_fashion_mnist} and \ref{fig:main_results_all_datasets}, listing average accuracy and average forgetting per training regime for each algorithm across CIFAR-100, MNIST, Fashion-MNIST, QMNIST, and KMNIST, together with standard deviations computed over task orderings.



\begin{table*}[t]
\centering
\caption{Average Accuracy and Average Forgetting per ordering $\pm$ std. across different training regimes on the CIFAR-100 dataset}
\resizebox{\textwidth}{!}{
\begin{tabular}{lcc cc cc cc cc}
\toprule
\multirow{2}{*}{\textbf{Algorithm / Regime}} 
& \multicolumn{2}{c}{\textbf{Last Block}} 
& \multicolumn{2}{c}{\textbf{Last 2 Blocks}} 
& \multicolumn{2}{c}{\textbf{Last 3 Blocks}} 
& \multicolumn{2}{c}{\textbf{Last 6 Blocks}} 
& \multicolumn{2}{c}{\textbf{Full Finetune (8 Blocks)}} \\
\cmidrule(lr){2-3} \cmidrule(lr){4-5} \cmidrule(lr){6-7} \cmidrule(lr){8-9} \cmidrule(lr){10-11}
& Acc & Forget 
& Acc & Forget 
& Acc & Forget 
& Acc & Forget 
& Acc & Forget \\
\midrule

EWC 
& $0.52 \pm 0.009$ & $0.02 \pm 0.008$ 
& $0.59 \pm 0.010$ & $0.02 \pm 0.006$ 
& $0.62 \pm 0.013$ & $0.02 \pm 0.009$ 
& $0.64 \pm 0.017$ & $0.03 \pm 0.006$ 
& $0.65 \pm 0.016$ & $0.03 \pm 0.008$ \\

LwF 
& $0.62 \pm 0.008$ & $0.05 \pm 0.005$ 
& $0.70 \pm 0.006$ & $0.08 \pm 0.004$ 
& $0.75 \pm 0.004$ & $0.08 \pm 0.005$ 
& $0.79 \pm 0.008$ & $0.07 \pm 0.007$ 
& $0.80 \pm 0.008$ & $0.07 \pm 0.008$ \\

SI 
& $0.57 \pm 0.007$ & $0.04 \pm 0.006$ 
& $0.61 \pm 0.013$ & $0.09 \pm 0.012$ 
& $0.62 \pm 0.021$ & $0.14 \pm 0.019$ 
& $0.62 \pm 0.018$ & $0.18 \pm 0.015$ 
& $0.63 \pm 0.017$ & $0.18 \pm 0.017$ \\

GEM 
& $0.52 \pm 0.008$ & $0.15 \pm 0.010$ 
& $0.55 \pm 0.010$ & $0.21 \pm 0.011$ 
& $0.57 \pm 0.012$ & $0.26 \pm 0.011$ 
& $0.61 \pm 0.012$ & $0.26 \pm 0.014$ 
& $0.71 \pm 0.008$ & $0.16 \pm 0.011$ \\

\bottomrule
\end{tabular}
}
\label{tab:cifar100}
\end{table*}


\begin{table*}[t]
\centering
\caption{Average Accuracy and Average Forgetting per ordering $\pm$ std. across different training regimes on the MNIST dataset}
\resizebox{\textwidth}{!}{
\begin{tabular}{lcc cc cc cc cc}
\toprule
\multirow{2}{*}{\textbf{Algorithm / Regime}} 
& \multicolumn{2}{c}{\textbf{Last Block}} 
& \multicolumn{2}{c}{\textbf{Last 2 Blocks}} 
& \multicolumn{2}{c}{\textbf{Last 3 Blocks}} 
& \multicolumn{2}{c}{\textbf{Last 6 Blocks}} 
& \multicolumn{2}{c}{\textbf{Full Finetune (8 Blocks)}} \\
\cmidrule(lr){2-3} \cmidrule(lr){4-5} \cmidrule(lr){6-7} \cmidrule(lr){8-9} \cmidrule(lr){10-11}
& Acc & Forget 
& Acc & Forget 
& Acc & Forget 
& Acc & Forget 
& Acc & Forget \\
\midrule

EWC 
& $0.98 \pm 0.012$ & $0.02 \pm 0.015$ 
& $0.91 \pm 0.082$ & $0.11 \pm 0.103$ 
& $0.87 \pm 0.083$ & $0.16 \pm 0.104$ 
& $0.86 \pm 0.060$ & $0.18 \pm 0.075$ 
& $0.89 \pm 0.078$ & $0.14 \pm 0.097$ \\

LwF 
& $0.99 \pm 0.002$ & $0.01 \pm 0.003$ 
& $0.99 \pm 0.007$ & $0.01 \pm 0.009$ 
& $0.99 \pm 0.007$ & $0.01 \pm 0.009$ 
& $0.99 \pm 0.004$ & $0.01 \pm 0.005$ 
& $0.99 \pm 0.011$ & $0.01 \pm 0.014$ \\

SI 
& $0.96 \pm 0.018$ & $0.04 \pm 0.021$ 
& $0.86 \pm 0.090$ & $0.17 \pm 0.113$ 
& $0.84 \pm 0.095$ & $0.20 \pm 0.119$ 
& $0.80 \pm 0.072$ & $0.25 \pm 0.090$ 
& $0.80 \pm 0.056$ & $0.25 \pm 0.070$ \\

GEM 
& $0.96 \pm 0.024$ & $0.04 \pm 0.030$ 
& $0.94 \pm 0.038$ & $0.08 \pm 0.047$ 
& $0.92 \pm 0.059$ & $0.10 \pm 0.075$ 
& $0.93 \pm 0.053$ & $0.09 \pm 0.067$ 
& $0.95 \pm 0.037$ & $0.06 \pm 0.047$ \\

\bottomrule
\end{tabular}
}
\label{tab:mnist}
\end{table*}



\begin{table*}[t]
\centering
\caption{Average Accuracy and Average Forgetting per ordering $\pm$ std. across different training regimes on the Fashion MNIST dataset}
\resizebox{\textwidth}{!}{
\begin{tabular}{lcc cc cc cc cc}
\toprule
\multirow{2}{*}{\textbf{Algorithm / Regime}} 
& \multicolumn{2}{c}{\textbf{Last Block}} 
& \multicolumn{2}{c}{\textbf{Last 2 Blocks}} 
& \multicolumn{2}{c}{\textbf{Last 3 Blocks}} 
& \multicolumn{2}{c}{\textbf{Last 6 Blocks}} 
& \multicolumn{2}{c}{\textbf{Full Finetune (8 Blocks)}} \\
\cmidrule(lr){2-3} \cmidrule(lr){4-5} \cmidrule(lr){6-7} \cmidrule(lr){8-9} \cmidrule(lr){10-11}
& Acc & Forget 
& Acc & Forget 
& Acc & Forget 
& Acc & Forget 
& Acc & Forget \\
\midrule

EWC 
& $0.92 \pm 0.043$ & $0.09 \pm 0.054$ 
& $0.91 \pm 0.057$ & $0.10 \pm 0.072$ 
& $0.86 \pm 0.063$ & $0.16 \pm 0.079$ 
& $0.81 \pm 0.082$ & $0.23 \pm 0.103$ 
& $0.85 \pm 0.041$ & $0.17 \pm 0.051$ \\

LwF 
& $0.91 \pm 0.042$ & $0.10 \pm 0.053$ 
& $0.90 \pm 0.059$ & $0.12 \pm 0.074$ 
& $0.86 \pm 0.084$ & $0.17 \pm 0.105$ 
& $0.85 \pm 0.077$ & $0.17 \pm 0.097$ 
& $0.82 \pm 0.049$ & $0.21 \pm 0.061$ \\

SI 
& $0.92 \pm 0.027$ & $0.09 \pm 0.035$ 
& $0.93 \pm 0.022$ & $0.08 \pm 0.027$ 
& $0.88 \pm 0.057$ & $0.14 \pm 0.071$ 
& $0.89 \pm 0.059$ & $0.13 \pm 0.075$ 
& $0.87 \pm 0.063$ & $0.16 \pm 0.080$ \\

GEM 
& $0.97 \pm 0.012$ & $0.02 \pm 0.013$ 
& $0.96 \pm 0.016$ & $0.03 \pm 0.019$ 
& $0.95 \pm 0.030$ & $0.05 \pm 0.041$ 
& $0.91 \pm 0.037$ & $0.08 \pm 0.051$ 
& $0.94 \pm 0.034$ & $0.06 \pm 0.042$ \\

\bottomrule
\end{tabular}
}
\label{tab:fashion}
\end{table*}

\begin{table*}[t]
\centering
\caption{Average Accuracy and Average Forgetting per ordering $\pm$ std. across different training regimes on the QMNIST dataset}
\resizebox{\textwidth}{!}{
\begin{tabular}{lcc cc cc cc cc}
\toprule
\multirow{2}{*}{\textbf{Algorithm / Regime}} 
& \multicolumn{2}{c}{\textbf{Last Block}} 
& \multicolumn{2}{c}{\textbf{Last 2 Blocks}} 
& \multicolumn{2}{c}{\textbf{Last 3 Blocks}} 
& \multicolumn{2}{c}{\textbf{Last 6 Blocks}} 
& \multicolumn{2}{c}{\textbf{Full Finetune (8 Blocks)}} \\
\cmidrule(lr){2-3} \cmidrule(lr){4-5} \cmidrule(lr){6-7} \cmidrule(lr){8-9} \cmidrule(lr){10-11}
& Acc & Forget 
& Acc & Forget 
& Acc & Forget 
& Acc & Forget 
& Acc & Forget \\
\midrule

EWC 
& $0.98 \pm 0.008$ & $0.01 \pm 0.010$ 
& $0.91 \pm 0.073$ & $0.10 \pm 0.090$ 
& $0.90 \pm 0.061$ & $0.12 \pm 0.076$ 
& $0.92 \pm 0.053$ & $0.10 \pm 0.067$ 
& $0.85 \pm 0.084$ & $0.19 \pm 0.104$ \\

LwF 
& $0.99 \pm 0.004$ & $0.01 \pm 0.005$ 
& $0.99 \pm 0.004$ & $0.01 \pm 0.004$ 
& $0.99 \pm 0.004$ & $0.01 \pm 0.005$ 
& $0.98 \pm 0.030$ & $0.03 \pm 0.036$ 
& $0.98 \pm 0.026$ & $0.03 \pm 0.033$ \\

SI 
& $0.97 \pm 0.025$ & $0.03 \pm 0.033$ 
& $0.85 \pm 0.095$ & $0.18 \pm 0.120$ 
& $0.79 \pm 0.127$ & $0.25 \pm 0.158$ 
& $0.82 \pm 0.068$ & $0.23 \pm 0.085$ 
& $0.85 \pm 0.101$ & $0.18 \pm 0.127$ \\

GEM 
& $0.96 \pm 0.014$ & $0.04 \pm 0.018$ 
& $0.93 \pm 0.047$ & $0.08 \pm 0.059$ 
& $0.93 \pm 0.045$ & $0.08 \pm 0.056$ 
& $0.94 \pm 0.052$ & $0.08 \pm 0.065$ 
& $0.95 \pm 0.049$ & $0.06 \pm 0.061$ \\

\bottomrule
\end{tabular}
}
\label{tab:qmnist}
\end{table*}



\begin{table*}[t]
\caption{Average Accuracy and Average Forgetting per ordering $\pm$ std. across different training regimes on the KMNIST dataset}
\centering
\resizebox{\textwidth}{!}{
\begin{tabular}{lcc cc cc cc cc}
\toprule
\multirow{2}{*}{\textbf{Algorithm / Regime}} 
& \multicolumn{2}{c}{\textbf{Last Block}} 
& \multicolumn{2}{c}{\textbf{Last 2 Blocks}} 
& \multicolumn{2}{c}{\textbf{Last 3 Blocks}} 
& \multicolumn{2}{c}{\textbf{Last 6 Blocks}} 
& \multicolumn{2}{c}{\textbf{Full Finetune (8 Blocks)}} \\
\cmidrule(lr){2-3} \cmidrule(lr){4-5} \cmidrule(lr){6-7} \cmidrule(lr){8-9} \cmidrule(lr){10-11}
& Acc & Forget 
& Acc & Forget 
& Acc & Forget 
& Acc & Forget 
& Acc & Forget \\
\midrule

EWC 
& $0.94 \pm 0.010$ & $-0.00 \pm 0.009$ 
& $0.95 \pm 0.016$ & $0.03 \pm 0.022$ 
& $0.93 \pm 0.028$ & $0.06 \pm 0.037$ 
& $0.92 \pm 0.033$ & $0.08 \pm 0.045$ 
& $0.93 \pm 0.026$ & $0.06 \pm 0.033$ \\

LwF 
& $0.95 \pm 0.009$ & $0.02 \pm 0.008$ 
& $0.97 \pm 0.012$ & $0.02 \pm 0.014$ 
& $0.97 \pm 0.008$ & $0.03 \pm 0.010$ 
& $0.96 \pm 0.008$ & $0.03 \pm 0.010$ 
& $0.97 \pm 0.009$ & $0.03 \pm 0.010$ \\

SI 
& $0.91 \pm 0.025$ & $0.05 \pm 0.037$ 
& $0.86 \pm 0.035$ & $0.14 \pm 0.045$ 
& $0.86 \pm 0.050$ & $0.15 \pm 0.061$ 
& $0.82 \pm 0.047$ & $0.20 \pm 0.057$ 
& $0.84 \pm 0.081$ & $0.17 \pm 0.099$ \\

GEM 
& $0.91 \pm 0.014$ & $0.07 \pm 0.017$ 
& $0.91 \pm 0.016$ & $0.10 \pm 0.019$ 
& $0.91 \pm 0.016$ & $0.10 \pm 0.021$ 
& $0.92 \pm 0.014$ & $0.09 \pm 0.018$ 
& $0.92 \pm 0.020$ & $0.09 \pm 0.026$ \\

\bottomrule
\end{tabular}
}
\label{tab:kmnist}
\end{table*}



\end{document}
